\ifdefined\gkver\else\def\gkver{student}\fi
\documentclass[\gkver]{gaokaozhenti}
\newcommand{\ccnum}[1]{{\fontspec{Noto Serif CJK SC}#1}}
\begin{document}

\chapter{2011年全国理综}
%% paperType: 理综物理部分

\begin{enumerate}
\item
%% number: 14
%% paperName: 2011年全国理综第 14 题 6 分
%% typeId: 2
%% type: 多选题
%% chapter: 气体性质
%% point: 热力学第一定律
%% method: 无
%% score: 6
%% degree: 650
%% duplicateId: 0
%% body:
关于一定量的气体，下列叙述正确的是\xzanswer{AD}
\fourchoices[answer=AD]
{气体吸收的热量可以完全转化为功}
{气体体积增大时，其内能一定减少}
{气体从外界吸收热量，其内能一定增加}
{外界对气体做功，气体内能可能减少}

%% answer: AD
%% memo:
\memoanswer{
由热力学第一定律 $Q + W = \Delta U$，在一定的条件下，控制温度不变，气体吸收的热量可以完全转化为功，其内能的变化为零，A 正确；气体体积增大时，对外界做功，但不知是否有热传递，无法确定内能的增加与减小，B 错误；气体从外界吸收热量，但不知气体是否对外界做功，无法确定内能的增加与减小，C 错误；外界对气体做功，但不知是否有热传递，无法确定内能的增加与减小，D 正确。
}

\item
%% number: 15
%% paperName: 2011年全国理综第 15 题 6 分
%% typeId: 1
%% type: 单选题
%% chapter: 磁场
%% point: 磁场的叠加
%% method: 无
%% score: 6
%% degree: 650
%% duplicateId: 0
%% body:
如图，两根相互平行的长直导线分别通有方向相反的电流$I_{1}$和$I_{2}$，且$I_{1}>I_{2}$；a、b、c、d 为导线某一横截面所在平面内的四点，且 a、b、c 与两导线共面；b 点在两导线之间，b、d 的连线与导线所在平面垂直。磁感应强度可能为零的点是\xzanswer{C}

\onepicture[w=5cm]{figs/15.png}

\fourchoices[answer=C]
{a 点}
{b 点}
{c 点}
{d 点}

%% answer: C
%% memo:
\memoanswer{
若某点合磁感应强度为零，必有 $I_{1}$ 和 $I_{2}$ 在该点形成磁场的磁感应强度等大反向，只有 c 点有可能，C 正确。
}

\item
%% number: 16
%% paperName: 2011年全国理综第 16 题 6 分
%% typeId: 1
%% type: 单选题
%% chapter: 光学
%% point: 光的折射
%% method: 无
%% score: 6
%% degree: 650
%% duplicateId: 0
%% body:
雨后太阳光入射到水滴中发生色散而形成彩虹。设水滴是球形的，图中的圆代表水滴过球心的截面，入射光线在过此截面的平面内，a、b、c、d 代表四条不同颜色的出射光线，则它们可能依次是\xzanswer{B}

\onepicture[h=4.5cm]{figs/16.png}

\fourchoices[answer=B]
{紫光、黄光、蓝光和红光}
{紫光、蓝光、黄光和红光}
{红光、蓝光、黄光和紫光}
{红光、黄光、蓝光和紫光}

%% answer: B
%% memo:
\memoanswer{
由题图可得：太阳光通过水滴折射后，偏折程度从小到大的排序为 d、c、b、a，则介质对色光的折射率为 $n_d < n_c < n_b < n_a$，色光的频率关系为 $f_d < f_c < f_b < f_a$，按题意，a、b、c、d 代表色光可能为紫光、蓝光、黄光和红光，B 正确。
}

\item
%% number: 17
%% paperName: 2011年全国理综第 17 题 6 分
%% typeId: 2
%% type: 多选题
%% chapter: 电场
%% point: 电势能电势差电场力做功
%% method: 无
%% score: 6
%% degree: 450
%% duplicateId: 0
%% body:
通常一次闪电过程历时约$0.2\sim0.3\Us$，它由若干个相继发生的闪击构成。每个闪击持续时间仅$40\sim80\Uus$，电荷转移主要发生在第一个闪击过程中。在某一次闪电前云地之间的电势差约为$1.0\times10^{9}\UV$，云地间距离约为$1\Ukm$；第一个闪击过程中云地间转移的电荷量约为$6\UC$，闪击持续时间约为$60\Uus$。假定闪电前云地间的电场是均匀的。根据以上数据，下列判断正确的是\xzanswer{AC}

\fourchoices[answer=AC]
{闪电电流的瞬时值可达到$1\times10^{5}\UA$}
{整个闪电过程的平均功率约为$1\times10^{14}\UW$}
{闪电前云地间的电场强度约为$1\times10^{6}\UVm$}
{整个闪电过程向外释放的能量约为$6\times10^{6}\UJ$}

%% answer: AC
%% memo:
\memoanswer{
第一个闪击过程的平均电流约为 $\overline{I} = \frac{q}{t} = \frac{6}{60 \times 10^{-6}}\UA = 1 \times 10^{5}\UA$，A 正确；第一个闪击过程的平均功率约为 $\overline{P} = U \overline{I} = 1.0 \times 10^{9} \times 1 \times 10^{5}\UW = 1 \times 10^{14}\UW$，但不等于整个闪电过程的平均功率，B 错误；$E = \frac{U}{d} = \frac{1.0 \times 10^{9}}{1 \times 10^{3}}\UVm = 1 \times 10^{6}\UVm$，C 正确；$W = q U = 6 \times 1.0 \times 10^{9} = 6 \times 10^{9}\UJ$，D 错误。
}

\item
%% number: 18
%% paperName: 2011年全国理综第 18 题 6 分
%% typeId: 1
%% type: 单选题
%% chapter: 原子和原子核
%% point: 玻尔模型
%% method: 无
%% score: 6
%% degree: 650
%% duplicateId: 0
%% body:
已知氢原子的基态能量为$E_{1}$，激发态能量$E_{n}=E_{1}/n^{2}$，其中$n=2,3,\cdots$。用$h$表示普朗克常量，$c$表示真空中的光速。能使氢原子从第一激发态电离的光子的最大波长为\xzanswer{C}

\fourchoices[answer=C]
{$-\frac{4hc}{3E_{1}}$}
{$-\frac{2hc}{3E_{1}}$}
{$-\frac{4hc}{E_{1}}$}
{$-\frac{9hc}{E_{1}}$}

%% answer: C
%% memo:
\memoanswer{
由氢原子的跃迁 $h\nu = E_m - E_n$ 及 $\nu = \frac{\lambda}{c}$，使氢原子从第一激发态电离，且第一激发态 $n = 2$，即氢原子从 $n = 2$ 跃迁到 $+\infty$，则有：$\frac{hc}{\lambda} = E_{+\infty} - E_2 = -\frac{E_1}{4}$，得：$\lambda = -\frac{4hc}{E_1}$，C 正确。
}

\item
%% number: 19
%% paperName: 2011年全国理综第 19 题 6 分
%% typeId: 1
%% type: 单选题
%% chapter: 曲线运动
%% point: 人造地球卫星
%% method: 无
%% score: 6
%% degree: 450
%% duplicateId: 0
%% body:
我国“嫦娥一号”探月卫星发射后，先在“24小时轨道”上绕地球运行（即绕地球一圈需要24小时）；然后，经过两次变轨依次到达“48小时轨道”和“72小时轨道”；最后奔向月球。如果按圆形轨道计算，并忽略卫星质量的变化，则在每次变轨完成后与变轨前相比，\xzanswer{D}

\fourchoices[answer=D]
{卫星动能增大，引力势能减小}
{卫星动能增大，引力势能增大}
{卫星动能减小，引力势能减小}
{卫星动能减小，引力势能增大}

%% answer: D
%% memo:
\memoanswer{
由 $G\frac{Mm}{r^{2}}=m\left(\frac{2\pi}{T}\right)^{2}r$ 得 $r=\sqrt[3]{\frac{GMT^{2}}{4\pi^{2}}}$，周期变长，轨道半径 $r$ 变大；又由 $G\frac{Mm}{r^{2}}=m\frac{v^{2}}{r}$ 得 $v=\sqrt{\frac{GM}{r}}$，轨道半径 $r$ 变大，速度减小，动能 $E_{k}=\frac{1}{2}mv^{2}$ 减小；轨道半径 $r$ 变大，引力做负功，引力势能增大，D 正确。
}

\item
%% number: 20
%% paperName: 2011年全国理综第 20 题 6 分
%% typeId: 2
%% type: 多选题
%% chapter: 运动和力
%% point: 动量和能量
%% method: 无
%% score: 6
%% degree: 450
%% duplicateId: 0
%% body:
质量为$M$、内壁间距为$L$的箱子静止于光滑的水平面上，箱子中间有一质量为$m$的小物块，小物块与箱子底板间的动摩擦因数为$\mu$。初始时小物块停在箱子正中间，如图所示。现给小物块一水平向右的初速度$v$，小物块与箱壁碰撞$N$次后恰又回到箱子正中间，并与箱子保持相对静止。设碰撞都是弹性的，则整个过程中，系统损失的动能为\xzanswer{BD}

\onepicture[w=5cm]{figs/20.png}

\fourchoices[answer=BD]
{$\frac{1}{2}mv^{2}$}
{$\frac{mM}{2(m+M)}v^{2}$}
{$\frac{1}{2}N\mu mgL$}
{$N\mu mgL$}

%% answer: BD
%% memo:
\memoanswer{
系统动量守恒，相当于完全非弹性碰撞，两物体最终速度相等设为 $u$，由动量守恒得 $mv = (m + M)u$，系统损失的动能为
$\frac{1}{2}mv^2 - \frac{1}{2}(m + M)u^2 = \frac{1}{2}\frac{mM}{m + M}v^2$，B 正确；碰撞 $N$ 次后恰又回到箱子正中间，小物块和箱子底板间的相对滑动距离 $d = NL$，而摩擦生热 $Q = fd = N\mu mgL$，系统损失的动能转化为内能，D 正确。
}

\item
%% number: 21
%% paperName: 2011年全国理综第 21 题 6 分
%% typeId: 1
%% type: 单选题
%% chapter: 机械波
%% point: 波长频率波速
%% method: 无
%% score: 6
%% degree: 450
%% duplicateId: 0
%% body:
一列简谐横波沿$x$轴传播，波长为$1.2\Um$，振幅为$A$。当坐标为$x=0$处质元的位移为$-\frac{\sqrt{3}}{2}A$向$y$轴负方向运动时，坐标为$x=0.4\Um$处质元的位移为$\frac{\sqrt{3}}{2}A$。当坐标为$x=0.2\Um$处的质元位于平衡位置且向$y$轴正方向运动时，$x=0.4\Um$处质元的位移和运动方向分别为\xzanswer{C}

\onepicture[w=6cm]{figs/21.png}

\fourchoices[answer=C]
{$-\frac{1}{2}A$，沿$y$轴正方向}
{$-\frac{1}{2}A$，沿$y$轴负方向}
{$-\frac{\sqrt{3}}{2}A$，沿$y$轴正方向}
{$-\frac{\sqrt{3}}{2}A$，沿$y$轴负方向}

%% answer: C
%% memo:
\memoanswer{
根据题意，此简谐横波必沿 $x$ 轴正向传播，在同一坐标系中做出两状态的波形图分别如实线和虚线所示，C 正确。
}

\item
%% number: 22
%% paperName: 2011年全国理综第 22 题 6 分
%% typeId: 4
%% type: 实验题
%% chapter: 气体性质
%% point: 用单分子油膜估测分子大小
%% method: 无
%% score: 6
%% degree: 650
%% duplicateId: 0
%% body:
在“油膜法估测油酸分子的大小”实验中，有下列实验步骤：

\begin{steps}
	\item
	往边长约为$40\Ucm$的浅盘里倒入约$2\Ucm$深的水，待水面稳定后将适量的痱子粉均匀地撒在水面上。
	\item
	用注射器将事先配好的油酸酒精溶液滴一滴在水面上，待薄膜形状稳定。
	\item
	将画有油膜形状的玻璃板平放在坐标纸上，计算出油膜的面积，根据油酸的体积和面积计算出油酸分子直径的大小。
	\item
	用注射器将事先配好的油酸酒精溶液一滴一滴地滴入量筒中，记下量筒内每增加一定体积时的滴数，由此计算出一滴油酸酒精溶液的体积。
	\item
	将玻璃板放在浅盘上，然后将油膜的形状用彩笔描绘在玻璃板上。
\end{steps}

完成下列填空：

\begin{enumerate}
	\item
	上述步骤中，正确的顺序是\tkanswer{④①②⑤③}（填写步骤前面的数字）。
	\item
	将$1\Ucmc$的油酸溶于酒精，制成$300\Ucmc$的油酸酒精溶液；测得$1\Ucmc$的油酸酒精溶液有50滴。现取一滴该油酸酒精溶液滴在水面上，测得所形成的油膜的面积是$0.13\Umq$。由此估算出油酸分子的直径为\tkanswer{$5\times10^{-10}$} $\Um$（结果保留1位有效数字）。
\end{enumerate}

%% answer: 
\jdanswer{
\begin{enumerate}
	\item
	④①②⑤③

	\item
	$5\times10^{-10}$

\end{enumerate}
}
%% memo:
\memoanswer{
（1）在“油膜法估测油酸分子的大小”实验中，按题目所给的实验步骤合理的顺序为④①②⑤③（填写步骤前面的数字）。

（2）1 滴油酸酒精溶液含油酸的体积 $V = \frac{1}{50} \times \frac{1}{300}\Ucmc = \frac{2}{3} \times 10^{-4}\Ucmc = \frac{2}{3} \times 10^{-10}\Umc$，油酸分子的直径约为 $d = \frac{V}{S} = \frac{\frac{2}{3} \times 10^{-10}}{0.13}\Um = 5 \times 10^{-10}\Um$。
}

\item
%% number: 23
%% paperName: 2011年全国理综第 23 题 12 分
%% typeId: 4
%% type: 实验题
%% chapter: 电路
%% point: 多用表的使用
%% method: 无
%% score: 12
%% degree: 450
%% duplicateId: 0
%% body:
使用多用电表测量电阻时，多用电表内部的电路可以等效为一个直流电源（一般为电池）、一个电阻和一表头相串联，两个表笔分别位于此串联电路的两端。现需要测量多用电表内电池的电动势，给定的器材有：待测多用电表，量程为$60\UmA$的电流表，电阻箱，导线若干。实验时，将多用电表调至$\times1\UO$挡，调好零点；电阻箱置于适当数值。完成下列填空：

\begin{enumerate}
	\item
	仪器连线如图\ref{2011全国理综23a}所示（a 和 b 是多用电表的两个表笔）。若两电表均正常工作，则表笔 a 为\tkanswer{黑}（填“红”或“黑”）色；
	\item
	若适当调节电阻箱后，图\ref{2011全国理综23a}中多用电表、电流表与电阻箱的示数分别如图\ref{2011全国理综23b}，\ref{2011全国理综23c}，\ref{2011全国理综23d}所示，则多用电表的读数为\tkanswer{14.0} $\UO$，电流表的读数为\tkanswer{53.0} $\UmA$，电阻箱的读数为\tkanswer{4.6} $\UO$；
	\item
	将图\ref{2011全国理综23a}中多用电表的两表笔短接，此时流过多用电表的电流为\tkanswer{102} $\UmA$；（保留3位有效数字）
	\item
	计算得到多用电表内电池的电动势为\tkanswer{1.54} $\UV$。（保留3位有效数字）
\end{enumerate}

\onepicture[num=1, label=2011全国理综23a, align=right, w=7cm]{figs/23a.png}
\threepicture[numA={2（a）}, numB={2（b）}, numC={2（c）}, labelA=2011全国理综23b, labelB=2011全国理综23c, labelC=2011全国理综23d, align=right, wA=4.6cm, wB=5.2cm, wC=3.2cm]
{figs/23b.png}{figs/23c.png}{figs/23d.png}

%% answer: 
\jdanswer{
\begin{enumerate}
	\item
	黑

	\item
	14.0，53.0，4.6

	\item
	102

	\item
	1.54

\end{enumerate}
}
%% memo:
\memoanswer{
（1）多用电表内部的电源正极与黑色表笔连接，电流表正常工作时电流从正接线柱流入，由此可以判定表笔 a 接内电源正极，为黑色。

（2）如图 2（a）为多用电表的$\times1\UO$欧姆挡，读数为 $14.0\UO$；如图 2（b）为电流表，读数为 $53.0\UmA$；如图 2（c）为电阻箱，读数为 $(0 \times 100 + 0 \times 10 + 4 \times 1 + 6 \times 0.1)\UO = 4.6\UO$。

（3）图 1 连接时，电路中电流 $I = 53.0\UmA$，则多用电表中电流为 $53.0\UmA$，指针指在电流格数为 26，将多用电表的两表笔短接，多用电表满偏为 50 格，则有：$\frac{50}{26} \times 53.0\UmA = 101.9\UmA = 102\UmA$。

（4）设多用电表在$\times1\UO$欧姆挡的内阻为 $r$，由于两表笔短接电流为 $102\UmA$，则 $E=102\times10^{-3}r$；图 1 连接时，$E=53.0\times10^{-3}\times(14.0+r)$，联立解得 $E=1.54\UV$，$r=15.1\UO$。
}

\item
%% number: 24
%% paperName: 2011年全国理综第 24 题 15 分
%% typeId: 6
%% type: 计算题
%% chapter: 电磁感应
%% point: 电磁感应电路分析
%% method: 无
%% score: 15
%% degree: 450
%% duplicateId: 0
%% body:
如图，两根足够长的金属导轨 ab、cd 竖直放置，导轨间距离为$L$，电阻不计。在导轨上端并接两个额定功率均为$P$、电阻均为$R$的小灯泡。整个系统置于匀强磁场中，磁感应强度方向与导轨所在平面垂直。现将一质量为$m$、电阻可以忽略的金属棒$MN$从图示位置由静止开始释放。金属棒下落过程中保持水平，且与导轨接触良好。已知某时刻后两灯泡保持正常发光。重力加速度为$g$。求：

\begin{enumerate}
	\item
	磁感应强度的大小；
	\item
	灯泡正常发光时导体棒的运动速率。
\end{enumerate}

\onepicture[align=right, w=4.5cm]{figs/24.png}

%% answer: 
\jdanswer{
\begin{enumerate}
	\item
	$B=\frac{mg}{2L}\sqrt{\frac{R}{P}}$

	\item
	$v=\frac{2P}{mg}$

\end{enumerate}
}
%% memo:
\memoanswer{
（1）设小灯泡的额定电流为 $I_{0}$，有

$P=I_{0}^{2}R$ ①

由题意，在金属棒沿导轨竖直下落的某时刻，小灯泡保持正常发光，流经 MN 电流为 $I=2I_{0}$ ②

此时金属棒 MN 所受的重力和安培力相等，下落的速度达到最大值，有

$mg=BIL$ ③

联立①②③式得

$B=\frac{mg}{2L}\sqrt{\frac{R}{P}}$ ④

（2）设灯泡正常发光时，导体棒的速率为 $v$，由电磁感应定律和欧姆定律得

$E=BLv$ ⑤

$E=RI_{0}$ ⑥

联立①②④⑤⑥式得

$v=\frac{2P}{mg}$ ⑦
}

\item
%% number: 25
%% paperName: 2011年全国理综第 25 题 19 分
%% typeId: 6
%% type: 计算题
%% chapter: 磁场
%% point: 带电粒子在组合场中的运动
%% method: 无
%% score: 19
%% degree: 450
%% duplicateId: 0
%% body:
如图，与水平面成$45^\circ$角的平面$MN$将空间分成Ⅰ和Ⅱ两个区域。一质量为$m$、电荷量为$q$（$q>0$）的粒子以速度$v_{0}$从平面$MN$上的$P_{0}$点水平向右射入Ⅰ区。粒子在Ⅰ区运动时，只受到大小不变、方向竖直向下的电场作用，电场强度大小为$E$；在Ⅱ区运动时，只受到匀强磁场的作用，磁感应强度大小为$B$，方向垂直于纸面向里。求粒子首次从Ⅱ区离开时到出发点$P_{0}$的距离。粒子的重力可以忽略。

\onepicture[align=right, w=5.2cm]{figs/25.png}

%% answer: 
\jdanswer{
$s=\frac{\sqrt{2}mv_{0}}{q}\left(\frac{2v_{0}}{E}+\frac{1}{B}\right)$
}
%% memo:
\memoanswer{
带电粒子进入电场后，在电场力的作用下沿抛物线运动，其加速度方向竖直向下，设其大小为 $a$，由牛顿定律得

$qE=ma$ ①

设经过时间 $t_{0}$，粒子从平面 $MN$ 上的点 $P_{1}$ 进入磁场，由运动学公式和几何关系得

$v_{0}t_{0}=\frac{1}{2}at_{0}^{2}$ ②

粒子速度大小 $V_{1}$ 为

$V_{1}=\sqrt{v_{0}^{2}+(at_{0})^{2}}$ ③

设速度方向与竖直方向的夹角为 $\alpha$，则

$\tan\alpha=\frac{v_{0}}{at_{0}}$ ④

\onepicture[w=7cm]{figs/25_an.png}

此时粒子到出发点 $P_{0}$ 的距离为

$s_{0}=\sqrt{2}v_{0}t_{0}$ ⑤

此后，粒子进入磁场，在洛伦兹力作用下做匀速圆周运动，圆周半径为

$r_{1}=\frac{mV_{1}}{qB}$ ⑥

设粒子首次离开磁场的点为 $P_{2}$，弧 $\widehat{P_{1}P_{2}}$ 所张的圆心角为 $2\beta$，则 $P_{1}$ 到点 $P_{2}$ 的距离为

$s_{1}=2r_{1}\sin\beta$ ⑦

由几何关系得

$\alpha+\beta=45^\circ$ ⑧

联立①②③④⑥⑦⑧式得

$s_{1}=\sqrt{2}\frac{mv_{0}}{qE}$ ⑨

点 $P_{2}$ 与点 $P_{0}$ 相距

$l=s_{0}+s_{1}$ \ccnum{⑩}

联立①②⑤⑨⑩式得

$l=\frac{\sqrt{2}mv_{0}}{q}\left(\frac{2v_{0}}{E}+\frac{1}{B}\right)$ \ccnum{⑪}
}

\item
%% number: 26
%% paperName: 2011年全国理综第 26 题 20 分
%% typeId: 6
%% type: 计算题
%% chapter: 运动和力
%% point: 动量和能量
%% method: 无
%% score: 20
%% degree: 250
%% duplicateId: 0
%% body:
装甲车和战舰采用多层钢板比采用同样质量的单层钢板更能抵御穿甲弹的射击。通过对以下简化模型的计算可以粗略说明其原因。质量为$2m$、厚度为$2d$的钢板静止在水平光滑桌面上。质量为$m$的子弹以某一速度垂直射向该钢板，刚好能将钢板射穿。现把钢板分成厚度均为$d$、质量均为$m$的相同两块，间隔一段距离水平放置，如图所示。若子弹以相同的速度垂直射向第一块钢板，穿出后再射向第二块钢板，求子弹射入第二块钢板的深度。设子弹在钢板中受到的阻力为恒力，且两块钢板不会发生碰撞。不计重力影响。

\twopicture[labelA=2011全国理综26a, labelB=2011全国理综26b, align=right, wA=6.5cm, wB=6.5cm]
{figs/26a.png}{figs/26b.png}

%% answer: 
\jdanswer{
$x=\frac{1}{2}\left(1+\frac{\sqrt{3}}{2}\right)d$
}
%% memo:
\memoanswer{
设子弹初速度为 $v_{0}$，射入厚度为 $2d$ 的钢板后，最终钢板和子弹的共同速度为 $V$。由动量守恒得

$(2m+m)V=mv_{0}$ ①

解得 $V=\frac{1}{3}v_{0}$。此过程中动能的损失为

$\Delta E=\frac{1}{2}mv_{0}^{2}-\frac{1}{2}\times3mV^{2}$ ②

解得 $\Delta E=\frac{1}{3}mv_{0}^{3}$。

分成两块钢板后，设子弹穿过第一块钢板时两者的速度分别为 $v_{1}$ 和 $V_{1}$，由动量守恒得 $mv_{1}+mV_{1}=mv_{0}$ ③

因为子弹在钢板中受到的阻力为恒力，射穿第一块钢板的动能损失为 $\frac{\Delta E}{2}$，由能量守恒得

$\frac{1}{2}mv_{1}^{2}+\frac{1}{2}mV_{1}^{2}=\frac{1}{2}mv_{0}^{2}-\frac{\Delta E}{2}$ ④

联立①②③④式，且考虑到 $v_{1}$ 必须大于 $V_{1}$，得

$v_{1}=\left(\frac{1}{2}+\frac{\sqrt{3}}{6}\right)v_{0}$ ⑤

设子弹射入第二块钢板并留在其中后两者的共同速度为 $V_{2}$，由动量守恒得

$2mV_{2}=mv_{1}$ ⑥

损失的动能为

$\Delta E'=\frac{1}{2}mv_{1}^{2}-\frac{1}{2}\times2mV_{2}^{2}$ ⑦

联立①②③⑤⑥⑦式得

$\Delta E'=\frac{1}{2}\left(1+\frac{\sqrt{3}}{2}\right)\times\frac{\Delta E}{2}$ ⑧

因为子弹在钢板中受到的阻力为恒力，由⑧式可得，射入第二块钢板的深度 $x$ 为

$x=\frac{1}{2}\left(1+\frac{\sqrt{3}}{2}\right)d$ ⑨
}

\end{enumerate}

\end{document}
